\section{Appendix C: Exact calibration and same-model alternatives}

The calibrated alternatives are specified by the following maps and continuous sign families.

The lower-bound constructions use exact calibration to combine a homogeneous conditional log-odds coefficient with continuous nuisance functions. We first describe the local equations and their smooth neighborhoods, then state the support guarantees for the finite sign families in \cref{obj:def:continuous-calibrated-priors}. The calibration maps are defined in \cref{obj:def:calibration-maps}.

The covariance adjustment \leanref{sym:\mathfrak c}{\ensuremath{\mathfrak c(t,\xi,\upsilon)}} prescribes the odds ratio of a binary table with treatment margin \leanref{sym:\xi}{\ensuremath{\xi}} and outcome margin \leanref{sym:\upsilon}{\ensuremath{\upsilon}}. Its normalized extension \leanref{sym:\mathfrak B}{\ensuremath{\mathfrak B(t,\xi,\upsilon)}} extends the effect-normalized adjustment smoothly through zero effect. The mixed equation uses the signed propensity amplitude \leanref{sym:\eta}{\ensuremath{\eta}} and signed outcome amplitude \leanref{sym:\zeta}{\ensuremath{\zeta}}; the fair equation uses the signed control-risk amplitude \leanref{sym:\delta}{\ensuremath{\delta}}. In the latter equation, the residual \leanref{sym:\mathfrak H}{\ensuremath{\mathfrak H(t,\delta,\xi,u)}} smoothly extends the matching discrepancy normalized by the effect and squared amplitude. The within-cell coordinate \leanref{sym:u}{\ensuremath{u}} ranges over \([0,1]\).

The scalar implicit-function theorem supplies the local graph representation needed for these calibration equations. Its conditions concern an open parameter domain, differentiability, and a nonzero derivative with respect to the scalar being calibrated.

\begin{lemmav}{lem:scalar-implicit-function}[Scalar implicit function theorem]
Let \(d\) be a nonnegative integer, let \(F:\mathbb R^d\times\mathbb R\to\mathbb R\), and let \(D\subseteq\mathbb R^d\times\mathbb R\). Fix \(\mathbf b_0\in\mathbb R^d\) and \(\chi_0\in\mathbb R\). Suppose that:
\begin{itemize}
\item \textbf{(Open domain.)} \(D\) is open and \((\mathbf b_0,\chi_0)\in D\).
\item \textbf{(Smoothness.)} \(F\) is continuously differentiable on \(D\).
\item \textbf{(Zero.)} \(F(\mathbf b_0,\chi_0)=0\).
\item \textbf{(Nonzero partial derivative.)} \(\partial_\chi F(\mathbf b_0,\chi_0)\ne0\).
\end{itemize}
Then there exist open sets \(U\subseteq\mathbb R^d\) and \(V\subseteq\mathbb R\), and a function \(q:\mathbb R^d\to\mathbb R\), such that
\[
\mathbf b_0\in U,\qquad \chi_0\in V,\qquad U\times V\subseteq D.
\]
The function \(q\) is continuously differentiable on \(U\), satisfies \(q(\mathbf b_0)=\chi_0\), and has \(q(\mathbf b)\in V\) for every \(\mathbf b\in U\). Moreover, for every \(\mathbf b\in U\) and every \(\chi\in V\),
\[
F(\mathbf b,\chi)=0
\quad\Longleftrightarrow\quad
\chi=q(\mathbf b).
\]
\end{lemmav}

The open-domain condition makes the calibrated scalar a differentiable function of the remaining parameters, including at parameter values that will lie on the boundary of the eventual support region. The local uniqueness conclusion identifies the zero set with that function throughout the product neighborhood. For the two calibration equations, \cref{obj:lem:exact-calibrations} supplies fixed brackets and common open neighborhoods for the resulting branches.

Specifically, the mixed neighborhood \(U\) contains a closed region with signed amplitude radius \(\varepsilon_{\mathrm{mix}}\), and the fair neighborhood \(V\) contains a closed region with signed amplitude radius \(\varepsilon_{\mathrm{fair}}\) and effect coordinate in \([0,1/4]\). Both amplitude radii exceed the calibration radius \(\varepsilon\) used for support. A further neighborhood \(W\) allows the fair margin to vary freely within a centered bracket of half-width \(b\); here \(b\) denotes a root-bracket width. The next statement establishes these neighborhoods together with the exact matching identities and derivative bounds.

\begin{lemmav}{lem:exact-calibrations}[Exact smooth calibration branches]
There exist constants \(\varepsilon>0\) and \(M>0\) with the following properties. Expectations below are over two independent fair signs, and
\[
Z_u=s_1\cos(\pi u/2)+s_2\sin(\pi u/2),
\qquad
E F(s_1,s_2)=\frac14\sum_{s_1,s_2\in\{-1,1\}}F(s_1,s_2).
\]
The covariance adjustment \(\mathfrak c\) and its normalized extension \(\mathfrak B\) are given by
\[
\begin{gathered}
d=e^t-1,\qquad
v_*=\xi(1-\xi)\upsilon(1-\upsilon),\qquad
L_*=1+d\{\xi(1-\upsilon)+(1-\xi)\upsilon\},\\
D(t)=\int_0^1 e^{rt}\,dr,\\
\mathfrak c(t,\xi,\upsilon)
=\frac{2d v_*}{L_*+\sqrt{L_*^2-4d^2v_*}},
\qquad
\mathfrak B(t,\xi,\upsilon)
=\frac{2D(t)v_*}{L_*+\sqrt{L_*^2-4d^2v_*}}.
\end{gathered}
\]
The mixed root \(\mathfrak q(\eta,\zeta,u)\) is selected from the roots in \((3/4,5/4)\) of
\[
16E\mathfrak B
 \bigl(32\eta\zeta,\tfrac12-\eta qZ_u,\tfrac12+\zeta Z_u\bigr)-q=0,
\]
with value \(1\) when that bracket contains no root. The fair root
\(\mathfrak p(t,\delta,u)\) is selected from the roots in \((3/10,1/2)\) of
\(\mathfrak H(t,\delta,p,u)=0\), with value \(2/5\) when that bracket contains no root. Here \(\mathfrak H\) is the smoothly extended normalized fair-matching residual, whose defining identity away from the axes is displayed below.

\begin{itemize}
\item \textbf{(Positive odds tables.)}
For every \(t,\xi,\upsilon\in\mathbb R\) satisfying
\[
|t|\le\varepsilon,\qquad
|\xi-\tfrac12|\le\varepsilon,\qquad
|\upsilon-\tfrac12|\le\varepsilon,
\]
put \(c=\mathfrak c(t,\xi,\upsilon)\). The four cells
\[
\begin{aligned}
p_{11}&=\xi\upsilon+c,&
p_{10}&=\xi(1-\upsilon)-c,\\
p_{01}&=(1-\xi)\upsilon-c,&
p_{00}&=(1-\xi)(1-\upsilon)+c
\end{aligned}
\]
are strictly positive and satisfy
\[
p_{11}+p_{10}=\xi,\qquad
p_{11}+p_{01}=\upsilon,\qquad
p_{11}+p_{10}+p_{01}+p_{00}=1,\qquad
\frac{p_{11}p_{00}}{p_{10}p_{01}}=e^t.
\]
Moreover,
\[
\mathfrak B(t,\xi,\upsilon)=\frac{\mathfrak c(t,\xi,\upsilon)}{t}
\quad(t\ne0),\qquad
\mathfrak B(0,\xi,\upsilon)=\xi(1-\xi)\upsilon(1-\upsilon).
\]

\item \textbf{(Signed mixed calibration.)}
For every \(|\eta|\le\varepsilon\), \(|\zeta|\le\varepsilon\), and \(u\in[0,1]\), put
\(q=\mathfrak q(\eta,\zeta,u)\). Then
\[
\frac34<q<\frac54,\qquad
16E\mathfrak B
 \bigl(32\eta\zeta,\tfrac12-\eta qZ_u,\tfrac12+\zeta Z_u\bigr)-q=0,
\]
and
\[
\begin{gathered}
\mathfrak q(\eta,\zeta,0)=\mathfrak q(\eta,\zeta,1),\\
E\mathfrak c
 \bigl(32\eta\zeta,\tfrac12-\eta qZ_u,\tfrac12+\zeta Z_u\bigr)
 =2\eta\zeta q,\\
\mathfrak q(-\eta,\zeta,u)=q,\qquad
\mathfrak q(\eta,-\zeta,u)=q,\\
\mathfrak q(0,\zeta,u)=1-4\zeta^2,\qquad
\mathfrak q(\eta,0,u)+4\eta^2\mathfrak q(\eta,0,u)^2=1.
\end{gathered}
\]

\item \textbf{(Fair calibration and separation.)}
For every \(t\in[0,1/4]\), \(|\delta|\le\varepsilon\), and \(u\in[0,1]\), put
\(p=\mathfrak p(t,\delta,u)\). Then
\[
\begin{gathered}
\mathfrak p(t,\delta,0)=\mathfrak p(t,\delta,1)=\frac25,\qquad
\mathfrak p(t,0,u)=\mathfrak p(0,\delta,u)=\frac25,\\
\mathfrak p(t,-\delta,u)=p,\qquad
\mathfrak H(t,\delta,p,u)=0,\\
E\mathsf G(t,p+\delta Z_u)=\mathsf G(\mathsf T(t,\delta),p).
\end{gathered}
\]
For \(t\delta\ne0\), the residual at the selected root agrees with the unnormalized fair-matching numerator divided by \(t\delta^2\). Furthermore,
\[
|p-\tfrac25|+|\partial_u\mathfrak p(t,\delta,u)|\le M\delta^2,
\]
and
\[
3t\delta^2\le t-\mathsf T(t,\delta)\le6t\delta^2,\qquad
\frac t2\le\mathsf T(t,\delta)\le t.
\]

\item \textbf{(Smooth neighborhoods and derivative bounds.)}
Both roots and every fully substituted local four-cell formula of the mixed and fair families are \(C^\infty\) on their respective closed parameter regions
\[
\begin{gathered}
|\eta|\le\varepsilon,\quad |\zeta|\le\varepsilon,\quad u\in[0,1],\\
t\in[0,1/4],\quad |\delta|\le\varepsilon,\quad u\in[0,1].
\end{gathered}
\]
For every integer \(m\ge0\), there is a constant \(B>0\) such that the norm of the \(m\)th total derivative of each root and each fully substituted local four-cell formula is at most \(B\) throughout its respective closed parameter region, uniformly over all family indices, auxiliary signs, and binary cell indices.

There also exist signed amplitude radii
\(\varepsilon_{\mathrm{mix}}>\varepsilon\) and
\(\varepsilon_{\mathrm{fair}}>\varepsilon\), a fixed fair-root bracket half-width
\(0<b<1/10\), and open sets \(U,V\subset\mathbb R^3\) containing, respectively,
\[
\begin{gathered}
\{(\eta,\zeta,u):|\eta|\le\varepsilon_{\mathrm{mix}},
                       \ |\zeta|\le\varepsilon_{\mathrm{mix}},\ u\in[0,1]\},\\
\{(t,\delta,u):t\in[0,1/4],\
                  |\delta|\le\varepsilon_{\mathrm{fair}},\ u\in[0,1]\}.
\end{gathered}
\]
Throughout \(U\), the selected mixed branch is the unique zero in its fixed bracket \((3/4,5/4)\); throughout \(V\), the selected fair branch is the unique zero in the fixed centered bracket \((2/5-b,2/5+b)\). Both roots and every fully substituted local four-cell formula are \(C^\infty\) on the corresponding open set.

Finally, there is an open set \(W\subset\mathbb R^4\) containing
\[
\{(t,\delta,\xi,u):t\in[0,1/4],\
 |\delta|\le\varepsilon_{\mathrm{fair}},\
 \xi\in(2/5-b,2/5+b),\ u\in[0,1]\}
\]
on which \(\mathfrak H\) is \(C^\infty\). Throughout this displayed core, for every margin \(\xi\) in the centered bracket and whenever \(t\delta\ne0\),
\[
\mathfrak H(t,\delta,\xi,u)
=\frac{E\mathsf G(t,\xi+\delta Z_u)-\mathsf G(\mathsf T(t,\delta),\xi)}
       {t\delta^2}.
\]
\end{itemize}
\end{lemmav}

The mixed calibration fixes the mean covariance adjustment exactly while allowing both margins to vary with the signs. The fair calibration instead matches the mean shifted risk, expressed through the logistic risk shift \leanref{sym:\mathsf G}{\ensuremath{\mathsf G(t,\xi)}} from \cref{obj:def:calibration-maps}, to a deterministic comparator. Its effect separation is between \(3t\delta^2\) and \(6t\delta^2\), while the comparator effect remains between \(t/2\) and \(t\). These identities specify the separation and matching simultaneously on the stated parameter regions.

The smoothness conclusion concerns the fully substituted cell probabilities as well as the calibration roots. Thus the derivative bounds incorporate the dependence of each cell on its calibrated margin. The open neighborhoods also include zero amplitudes and zero effect, where the normalized equations have their smooth extensions. In the fair family, both the displacement of the calibrated margin from \(2/5\) and its within-cell derivative are bounded by a constant times the squared amplitude.

We now apply these local maps to the shared-endpoint construction in \cref{obj:def:continuous-calibrated-priors}. A finite endpoint sign vector \leanref{sym:\sigma}{\ensuremath{\sigma}} determines the trigonometric field on the equal partition, with adjacent cells sharing the sign at their common endpoint. The endpoint identities in \cref{obj:lem:exact-calibrations} make the calibrated margins compatible with this continuous field. The following statement certifies model membership for every sign realization, together with the exact singleton matching and uniform amplitude derivatives needed for the sample mixtures.

\begin{lemmav}{lem:calibrated-support}[Uniform calibrated support and smoothness]
There exist jointly chosen absolute constants \(\varepsilon\), the calibration radius, and \(M\), the derivative bound, satisfying
\[
0<\varepsilon\le \frac{1}{100},
\qquad M\ge 1,
\]
with the following properties for the calibrated families of \cref{obj:def:continuous-calibrated-priors}. Throughout, \(k\ge1\) is an integer and \(\sigma\in\{-1,1\}^{k+1}\) is a sign vector.

\begin{itemize}
\item \textbf{(Mixed membership and effects.)}
For every \((\alpha,\beta)\in\mathcal D\), the prescribed domain of public smoothness exponents, set
\[
\eta=\varepsilon k^{-\alpha},
\qquad
\zeta=\varepsilon k^{-\beta}.
\]
The calibrated mixed laws \(P_{0,\sigma}\) and \(P_{1,\sigma}\) at these amplitudes belong to \(\mathcal M(\alpha,\beta)\) as defined in \cref{obj:def:model}, and their homogeneous conditional log-odds coefficients satisfy
\[
\theta(P_{0,\sigma})=0,
\qquad
\theta(P_{1,\sigma})=32\eta\zeta.
\]
For each treatment label, outcome label, and covariate \(x\), the sums over all sign vectors of the corresponding conditional cell probabilities under the two mixed laws are equal.

\item \textbf{(Fair membership and effects.)}
For every \((\alpha,\beta)\in\mathcal D\), set
\[
\delta=\varepsilon k^{-\beta}.
\]
For every \(t\in[0,1/4]\), the fair random law \(P_{t,\sigma}\) and the deterministic fair comparator \(P_{*,t,\delta}\) belong to \(\mathcal M(\alpha,\beta)\), and
\[
\theta(P_{t,\sigma})=t,
\qquad
\theta(P_{*,t,\delta})=\mathsf T(t,\delta),
\]
where \(\mathsf T(t,\delta)\) is the calibrated comparator effect defined in \cref{obj:def:calibration-maps}. For each treatment label, outcome label, and covariate \(x\), the uniform average over the \(2^{k+1}\) sign vectors of the conditional cell probability under \(P_{t,\sigma}\) equals the corresponding comparator cell probability.

\item \textbf{(Signed validity and derivative bounds.)}
For all signed amplitudes satisfying
\[
|\eta|\le\varepsilon,
\qquad
|\zeta|\le\varepsilon,
\qquad
|\delta|\le\varepsilon,
\]
the conditional four-cell tables of both mixed branches are valid probability tables, as are those of both fair branches for every \(t\in[0,1/4]\). For every sign vector, branch, treatment label, outcome label, and covariate \(x\), their one-record densities relative to uniform four-label measure lie in \([1/2,2]\). The mixed densities, viewed as functions of \((\eta,\zeta)\), and the fair densities, viewed as functions of \(\delta\) with \(t\) fixed, are twice continuously differentiable at every point of the corresponding closed amplitude box. The norms of their first and second total amplitude derivatives are at most \(M\), uniformly over all these indices and amplitudes. These bounds include the mixed second derivative in \(\eta\) and \(\zeta\) and the deterministic comparator's derivatives in \(\delta\).

\item \textbf{(Common smooth neighborhoods.)}
There exists an open subset of \(\mathbb R^2\) containing
\[
\{(\eta,\zeta):|\eta|\le\varepsilon,\ |\zeta|\le\varepsilon\}
\]
on which every mixed density is infinitely continuously differentiable in its amplitude arguments. There also exists an open subset of \(\mathbb R\) containing
\[
[-\varepsilon,\varepsilon]
\]
on which every fair density is infinitely continuously differentiable in \(\delta\). Each neighborhood is common to every integer \(k\ge1\), sign vector, branch, treatment label, outcome label, and covariate \(x\); the fair neighborhood is also common to every \(t\in[0,1/4]\).

\item \textbf{(Signed singleton matching.)}
For every \(|\eta|\le\varepsilon\) and \(|\zeta|\le\varepsilon\), and each treatment label, outcome label, and covariate \(x\), the sums over all sign vectors of the corresponding conditional cell probabilities under the two mixed laws are equal. For every \(|\delta|\le\varepsilon\) and \(t\in[0,1/4]\), the uniform average over all \(2^{k+1}\) sign vectors of each conditional cell probability under the fair random law equals the corresponding cell probability under the deterministic fair comparator.
\end{itemize}
\end{lemmav}

Membership in \cref{obj:def:model} certifies the prescribed envelopes and Hölder restrictions for the native propensity logit \leanref{sym:g}{\ensuremath{g(P,\cdot)}} and prognosis logit \(\nu(P,\cdot)\), together with the homogeneous logistic relation. This certification applies to every sign realization at the stated amplitude scales and throughout the prescribed domain of public exponents. Continuity across partition endpoints and the homogeneous scalar coefficient therefore hold for each law entering the finite priors in \cref{obj:def:continuous-calibrated-priors}.

Exact singleton matching concerns every conditional cell probability at every covariate value. For the mixed families, the two uniform sign averages agree despite their distinct homogeneous coefficients. For the fair family, the uniform sign average agrees with the deterministic comparator at its calibrated coefficient. The matching identities continue to hold throughout the signed amplitude boxes, providing exact cancellations alongside the first and second amplitude derivative bounds.

Finally, the density bounds and common smooth neighborhoods in \cref{obj:lem:calibrated-support} are uniform in the partition size, sign vector, labels, and covariate. They furnish a joint regularity certificate for the actual one-record densities, including the mixed cross derivative and the comparator's amplitude derivatives. Together with the effect identities, these properties provide same-model alternatives with controlled separation for the original-record lower-bound argument.
